\documentclass{article}
\usepackage[utf8]{inputenc}
\usepackage{amsmath,amssymb}
\usepackage[most]{tcolorbox}
\usepackage{geometry}
\geometry{margin=2.5cm}

\newcommand{\step}[1]{\par\noindent\hangindent=3.5em\hangafter=1%
  \makebox[3.5em][l]{\textbf{#1.}}}
\newcommand{\steptag}[1]{\hfill{\small\textsf{[#1]}}}

\newtcolorbox{proofbox}[1]{
  colback=gray!5, colframe=gray!60,
  fonttitle=\bfseries, title={#1},
  breakable, enhanced,
  left=4pt, right=4pt, top=4pt, bottom=4pt,
  before skip=10pt, after skip=10pt
}

\begin{document}

\begin{proofbox}{EXT-CB: there are universal C\_ext < infinity and e\_ext > ...}
\step{1} EXT-CB: there are universal C\_ext < infinity and e\_ext > 0 such that if e=delta+epsilon <= e\_ext, P,Q are delta-projections in an extended epsilon-C*-algebra A with ||P+Q-I|| <= delta, v:M\_r->S\_P is an extended delta-isomorphism, dim S\_Q=1 at level one, and S\_\{P,Q\} is nonzero, then there is one map v\_+:M\_\{r+1\}->A whose every amplification is a C\_ext*e-isomorphism; the same level-one unitary and the same four corner maps carry all amplification levels, with constants independent of r, n, and dim A. \steptag{validated} \\[6pt]
\step{1.1} Setup and level-one dimensions. Assume the root hypotheses and put e=delta+epsilon, P\_1=P, P\_2=Q, H\_j=S\_\{P\_j,Q\}, K\_1=C\textasciicircum{}r and K\_2=C. By def-extcb-datum these hypotheses are an EXT-CB datum. Choose e below the universal threshold e\_sel of lem-extcb1-cross-corner-dimension. That external gives dim H\_1=dim S\_\{P,Q\}=r and dim H\_2=dim S\_\{Q,Q\}=1. This is exclusively a level-one use of dim S\_Q=1; no amplification of Q is called one-dimensional. Choose once and for all a unitary U\_2:K\_2->H\_2. \steptag{validated} \\[4pt]
\step{1.2} Dimension-free exact-target correction lemma. There are universal a\_corr>0 and C\_corr<infinity with the following property: if B is a finite-dimensional C*-algebra, H a finite-dimensional Hilbert space, and T:B->B(H) is linear, dagger-preserving, has ||T\_n(XY)-T\_n(X)T\_n(Y)||<=a||X||||Y|| and ||T\_n(I)-I||<=a at every n, where 0<=a<=a\_corr, then one unital dagger-homomorphism mu:B->B(H) satisfies ||mu\_n-T\_n||<=C\_corr*a at every n. The same level-one mu is used at every amplification. This lemma is proved in children 1.2.1--1.2.3 from the norm-one diagonal and Newton correction, with no external theorem. \steptag{validated} \\[6pt]
\step{1.2.1} Uniform cb control, exact unitalization, and the norm-one diagonal. At each n and ||X||=1, dagger preservation and the defect estimate give ||T\_n(X)||\textasciicircum{}2<=||T\_n(X\textasciicircum{}dagger X)||+a, so M\_n=||T\_n|| satisfies M\_n\textasciicircum{}2<=M\_n+a and hence ||T||\_cb<=2 for a<=1. Fix a state phi on B and define S\_0(x)=T(x)+phi(x)(I-T(I)). States have cb norm one and phi(x\textasciicircum{}dagger)=conj(phi(x)), so S\_0 is dagger-preserving, exactly unital, ||S\_0-T||\_cb<=a, ||S\_0||\_cb<=3, and direct expansion gives multiplicative defect b\_0<=c\_u*a with the numerical universal c\_u=7. For every finite-dimensional C*-algebra B, normalized Haar measure on its compact unitary group gives D=int U\textasciicircum{}dagger tensor U dU. In finite dimension D is a finite convex combination sum\_s p\_s U\_s\textasciicircum{}dagger tensor U\_s with projective norm one; Haar invariance gives XD=DX and multiplication(D)=I. For a defect G of S, w'(x)=sum\_s S(p\_s U\_s\textasciicircum{}dagger)G(U\_s,x) satisfies ||w'||\_cb<=||S||\_cb||G||. At level n the same formula is entrywise I\_n tensor w', because the diagonal identity moves arbitrary matrix entries. Thus every correction is one level-one map with bounds independent of n and all dimensions. \steptag{validated} \\[4pt]
\step{1.2.2} One normalized Newton correction. Let S:B->B(H) be exactly unital, dagger-preserving, ||S||\_cb<=4, with multiplicative defect G of uniform amplified bilinear norm b. Use node 1.2.1 and put w'(x)=sum\_s S(p\_s U\_s\textasciicircum{}dagger)G(U\_s,x), w''(x)=w'(x\textasciicircum{}dagger)\textasciicircum{}dagger, w=(w'+w'')/2, and S\textasciicircum{}+=S+w. Since G(y,I)=G(I,y)=0, w(I)=0; hence S\textasciicircum{}+ remains unital and dagger-preserving, with ||w||\_cb<=4b. Exact associativity gives S(x)G(y,z)-G(xy,z)+G(x,yz)-G(x,y)S(z)=0. Expanding S(x)w(y)-w(xy)+w(x)S(y), using XD=DX, multiplication(D)=I, and S(I)=I, cancels G exactly to first order; every remaining term contains two defect factors. A direct triangle estimate therefore gives ||G\_\{S\textasciicircum{}+\}||<=K\_N*b\textasciicircum{}2 at every amplification for one numerical universal K\_N, while ||S\textasciicircum{}+-S||\_cb<=4b. Node 1.2.1 makes the identities uniform and ensures S\textasciicircum{}+ is the amplification of the same corrected level-one map. \steptag{validated} \\[4pt]
\step{1.2.3} Convergence and exactness. Choose a\_corr>0 so K\_N*c\_u*a\_corr<=1/2 and 8*c\_u*a\_corr<1. Start from the exactly unital S\_0 of node 1.2.1 and apply node 1.2.2 recursively. With b\_0<=c\_u*a and b\_\{k+1\}<=K\_N*b\_k\textasciicircum{}2, one has b\_k<=b\_0*2\textasciicircum{}\{-k\}. Thus ||S\_k-S\_0||\_cb<=sum\_k 4b\_k<=8c\_u*a<1; since ||S\_0||\_cb<=3, every S\_k has cb norm below 4 and every iteration is justified. The sequence converges in cb norm to one unital dagger-preserving linear map mu, and passing to the limit gives mu(xy)=mu(x)mu(y). Moreover ||mu-T||\_cb<=||S\_0-T||\_cb+||mu-S\_0||\_cb<=(1+8c\_u)a, and amplification commutes with the limit. Hence node 1.2 holds with the universal C\_corr=1+8c\_u=57, using the same level-one mu at all n. \steptag{validated} \\[6pt]
\step{1.2.3.1} Validated-dependency and induction bridge. Nodes 1.2.1 and 1.2.2 are now validated. Let b\_k be the uniform amplified bilinear norm of the multiplicative defect of S\_k. Node 1.2.1 gives an exactly unital dagger-preserving S\_0 with ||S\_0||\_cb<3, b\_0<=c\_u*a, and ||S\_0-T||\_cb<=a. Assume inductively that S\_k is exactly unital and dagger-preserving, ||S\_k||\_cb<4, and b\_k<=b\_0*2\textasciicircum{}\{-k\}. Node 1.2.2 applies and produces the same level-one map S\_\{k+1\} at all amplifications, still exactly unital and dagger-preserving, with ||S\_\{k+1\}-S\_k||\_cb<=4b\_k and b\_\{k+1\}<=K\_N*b\_k\textasciicircum{}2. Since K\_N*b\_0<=K\_N*c\_u*a<=1/2, b\_\{k+1\}<=b\_k/2<=b\_0*2\textasciicircum{}\{-(k+1)\}. Moreover ||S\_k-S\_0||\_cb<=sum\_\{j<k\}4b\_j<=8b\_0<=8c\_u*a<1, hence ||S\_k||\_cb<4 and the induction closes. This estimate also covers b\_0=0: then b\_k=0 for every k and every correction increment vanishes. Therefore sum\_k ||S\_\{k+1\}-S\_k||\_cb<=8b\_0, so S\_k converges in cb norm to one exactly unital dagger-preserving map mu. For each x,y, continuity of multiplication and b\_k->0 give mu(xy)-mu(x)mu(y)=lim\_k [S\_k(xy)-S\_k(x)S\_k(y)]=0. Cb convergence means sup\_n ||(S\_k-mu)\_n|$|-\rangle$0, so every mu\_n is the amplification of this same level-one mu and ||mu-T||\_cb<=a+8c\_u*a=57a (c\_u=7). Thus the application of the now-validated Newton correction is permitted and all its hypotheses are verified at every iteration. \steptag{validated} \\[6pt]
\step{1.2.3.1.1} Degenerate-safe geometric estimate. Replace the strict bound in the parent node by ||S\_k-S\_0||\_cb <= sum\_\{j<k\} 4 b\_j <= 4 b\_0 sum\_\{j=0\}\textasciicircum{}\{k-1\}2\textasciicircum{}\{-j\} <= 8 b\_0 <= 8 c\_u a < 1. This remains valid when b\_0=0: then b\_k<=b\_0 2\textasciicircum{}\{-k\}=0 for every k, so every b\_k=0 and the correction bounds give ||S\_\{k+1\}-S\_k||\_cb<=4b\_k=0. In all cases, ||S\_k||\_cb <= ||S\_0||\_cb+||S\_k-S\_0||\_cb < 3+1=4, because ||S\_0||\_cb<3 and the displacement bound is strictly less than 1. Thus the induction closes without assuming b\_0>0. Likewise, sum\_\{k>=0\}||S\_\{k+1\}-S\_k||\_cb <= 4 sum\_\{k>=0\}b\_k <= 8b\_0; no strict geometric-series inequality is required. \steptag{validated} \\[4pt]
\step{1.3} Construction of one spatial four-corner system. Define h\_jk=Ha\textasciicircum{}Q\_\{P\_j,P\_k\} and T=h\_11 composed with v. The product, unit, upper-norm and adjoint clauses of conj-hcb together with def-extended-delta-inclusion imply, uniformly in n, that T is a dagger-preserving extended A\_0*e-homomorphism for a universal A\_0 (for example after enlarging C\_H, A\_0=4(C\_H+1)). Apply node 1.2 at e small enough to get one exact unital dagger-homomorphism mu\_11:M\_r->B(H\_1) with ||mu\_11,n-h\_11,n v\_n||<=kappa*e, kappa=C\_corr*A\_0. It is nonzero, hence injective because M\_r is simple, and it is onto because both spaces have dimension r\textasciicircum{}2 by node 1.1. Thus it is a spatial dagger-isomorphism: choose one unitary U\_1:C\textasciicircum{}r->H\_1 with mu\_11(A)=U\_1 A U\_1\textasciicircum{}dagger. With the already fixed U\_2, define mu\_jk(A)=U\_j A U\_k\textasciicircum{}dagger for A in B(K\_k,K\_j). These four fixed level-one maps are exact, completely isometric, adjoint-compatible and multiplicative at all amplifications via I\_n tensor U\_j; no level-dependent choice is made. \steptag{validated} \\[6pt]
\step{1.3.1} The amplified defect of T. For X,Y in M\_n tensor M\_r, def-extended-delta-inclusion gives ||v\_n(XY)-v\_n(X) dot v\_n(Y)||<=delta||X||||Y||, ||v\_n(X)||<=(1+delta)||X||, exact dagger preservation, and the corresponding unit estimate in the compressed algebra S\_P. The upper, product, unit and adjoint clauses of conj-hcb for h\_11,n therefore give ||T\_n(XY)-T\_n(X)T\_n(Y)|| <= (1+C\_H*e)delta||X||||Y||+C\_H*e*(1+delta)\textasciicircum{}2||X||||Y|| and ||T\_n(I)-I||<=A\_0*e; dagger is exact. For e below one universal bound, both defects are <=A\_0*e with A\_0=4(C\_H+1), independent of n,r and dim A. Thus node 1.2 applies and supplies one exact unital dagger-homomorphism mu\_11 with complete distance at most kappa*e, kappa=C\_corr*A\_0. \steptag{validated} \\[6pt]
\step{1.3.1.1} Endpoint/positive-parameter case split closing the correction step. The error tolerances delta and epsilon are nonnegative, so e=delta+epsilon is nonnegative. If e=0, then delta=epsilon=0. Substitution in the already established amplified estimates of node 1.3.1 gives, for every n and X,Y, ||T\_n(XY)-T\_n(X)T\_n(Y)||<=0 and ||T\_n(I)-I||<=0, while dagger preservation is exact. Hence the linear map T:M\_r->B(H\_1) is itself an exact unital dagger-homomorphism at level one (and its fixed amplifications have the same exact identities); define mu\_11:=T. Then ||mu\_11,n-T\_n||=0=kappa*e for every n. If e>0, shrink the universal threshold further so that e<a\_corr/A\_0 and retain the preceding defect bound with correction parameter a=A\_0*e. Then 0<a<a\_corr, node 1.2 applies, and its single level-one unital dagger-homomorphism mu\_11 satisfies ||mu\_11,n-T\_n||<C\_corr*A\_0*e=kappa*e for every n, hence also the asserted at-most bound. These two cases exhaust e>=0 and use no level-, rank-, or dimension-dependent choice. \steptag{validated} \\[6pt]
\step{1.3.1.1.1} Validated-dependency bridge for the correction step. Require node 1.2 to be validated before the positive-parameter conclusion of this child is available. Let A\_0=4(C\_H+1)>0 and kappa=C\_corr*A\_0, and shrink the universal EXT-CB threshold so e<=a\_corr/(2*A\_0). Node 1.3.1 gives, at every amplification, exact dagger preservation and multiplicative and unit defects at most A\_0*e for the same level-one map T. If e=0, nonnegativity of delta and epsilon gives delta=epsilon=0, so both defects vanish; hence T itself is an exact unital dagger-homomorphism and mu\_11:=T has ||mu\_11,n-T\_n||=0=kappa*e for every n, without using node 1.2. If e>0, set a=A\_0*e. Then 0<a<=a\_corr/2<a\_corr, so all hypotheses of the now-required validated node 1.2 hold. It supplies one level-one unital dagger-homomorphism mu\_11 whose fixed amplifications satisfy ||mu\_11,n-T\_n||<C\_corr*a=kappa*e, hence <=kappa*e, for every n. Thus the correction conclusion follows on the full e>=0 scope, but the e>0 branch is explicitly blocked until node 1.2 is validated; no unvalidated dependency is used. \steptag{archived} \\[4pt]
\step{1.3.1.1.2} Dependency-gated positive-parameter application. This child requires validated node 1.2. Under the universal threshold e<=a\_corr/(2*A\_0), the e>0 branch has correction parameter a=A\_0*e satisfying 0<a<a\_corr. The amplified defect and unit estimates of node 1.3.1 and exact dagger preservation verify every hypothesis of node 1.2 for the same level-one map T, so validated node 1.2 yields one unital dagger-homomorphism mu\_11 with ||mu\_11,n-T\_n||<C\_corr*A\_0*e=kappa*e at every n. The separate e=0 calculation in node 1.3.1.1 uses mu\_11=T and needs no correction lemma. Consequently the complete e>=0 conclusion is available only after node 1.2 is validated. \steptag{validated} \\[4pt]
\step{1.3.2} Spatiality and all exact corners. By node 1.1 dim H\_1=r. The unital mu\_11 is nonzero; its kernel is a two-sided ideal in the simple algebra M\_r, hence zero. Domain M\_r and codomain B(H\_1) both have dimension r\textasciicircum{}2, so mu\_11 is onto. A unital dagger-isomorphism between full matrix algebras preserves rank-one projections and matrix units, hence choosing an orthonormal basis in their ranges gives a single unitary U\_1:C\textasciicircum{}r->H\_1 with mu\_11(A)=U\_1 A U\_1\textasciicircum{}dagger. Together with the single U\_2 from node 1.1, mu\_jk(A)=U\_j A U\_k\textasciicircum{}dagger is a bijective linear isometry B(K\_k,K\_j)->B(H\_k,H\_j). Direct multiplication and adjoint calculation gives mu\_jl(XY)=mu\_jk(X)mu\_kl(Y), mu\_kj(X\textasciicircum{}dagger)=mu\_jk(X)\textasciicircum{}dagger and mu\_jj(I)=I. Tensoring the same identities with I\_n proves complete isometry and all amplified identities without new choices. \steptag{validated} \\[4pt]
\step{1.4} Conditional Ha inverses, in the required order. Shrink the universal threshold so kappa*e<1 and C\_H*e<1/4. Since mu\_11 is invertible and ||T-mu\_11||<1, T is bijective by Neumann inversion; v is bijective, so h\_11 is bijective. Moreover for Z=v(A), ||h\_11(Z)||/||Z|| >= (1-kappa*e)/(1+delta)>1/4. The canonical-identity closeness clause of conj-hcb makes h\_12,h\_21,h\_22 level-one bijections by Neumann inversion and gives h\_22 level-one lower modulus at least 1/4. Only now apply the conditional clauses of conj-hcb: the diagonal clause to h\_11 and h\_22, the off-diagonal clause to h\_12 anchored at h\_22, and to h\_21 anchored at h\_11. Hence every h\_jk,n is bijective, (1-C\_H*e)||Z||<=||h\_jk,n(Z)||<=(1+C\_H*e)||Z|| and ||h\_jk,n\textasciicircum{}\{-1\}||<=1+C\_H*e. Algebraically h\_jk,n\textasciicircum{}\{-1\}=I\_n tensor h\_jk\textasciicircum{}\{-1\}, so the inverses are not chosen afresh at level n. \steptag{validated} \\[6pt]
\step{1.4.1} Level-one triggers are established before inversion clauses. At n=1, identify T with mu\_11 followed by I+mu\_11\textasciicircum{}\{-1\}(T-mu\_11); since ||T-mu\_11||<=kappa*e<1 and mu\_11 is isometric, the parenthesized operator is invertible by its Neumann series. Thus T and then h\_11=T composed with v\textasciicircum{}\{-1\} are bijective. If Z=v(A), then ||h\_11(Z)||>=||mu\_11(A)||-kappa*e||A||=(1-kappa*e)||A||, whereas ||Z||<= (1+delta)||A||, so the lower modulus of h\_11 is at least (1-kappa*e)/(1+delta)>1/4. Separately, the canonical-identity closeness estimates in conj-hcb make each of h\_12,h\_21,h\_22 a strict-norm perturbation of its canonical level-one bijection when e is universally small. Neumann inversion gives their level-one bijectivity, and for h\_22 the same estimate gives lower modulus >1/4. \steptag{validated} \\[4pt]
\step{1.4.2} Apply precisely the conditional H-CB clauses. The established bijectivity and lower modulus allow the diagonal clause of conj-hcb first for h\_11 and h\_22, yielding bijectivity of every amplification and inverse norm <=1+C\_H*e. Then apply its off-diagonal clause to h\_12 with diagonal anchor h\_22 and to h\_21 with diagonal anchor h\_11; all hypotheses were established in node 1.4.1. Together with the unconditional upper bounds this gives the two-sided estimates in node 1.4 for every j,k,n. Finally, for a level-one linear bijection h\_jk, tensor algebra gives (I\_n tensor h\_jk)(I\_n tensor h\_jk\textasciicircum{}\{-1\})=I and the reverse product, so uniqueness of inverse gives h\_jk,n\textasciicircum{}\{-1\}=I\_n tensor h\_jk\textasciicircum{}\{-1\}. \steptag{validated} \\[4pt]
\step{1.5} Four fixed corner maps form a merging datum. Define gamma\_11=v and gamma\_jk=h\_jk\textasciicircum{}\{-1\} composed with mu\_jk for (j,k)!=(1,1). These are four bijective level-one maps and their amplifications are fixed. Put d\_jk,n=h\_jk,n gamma\_jk,n-mu\_jk,n. Then ||d\_11,n||<=kappa*e while d\_12,n=d\_21,n=d\_22,n=0. Exact adjointness from conj-hcb and the spatial system gives gamma\_kj,n(X\textasciicircum{}dagger)=gamma\_jk,n(X)\textasciicircum{}dagger. Using the product-defect and inverse bounds from conj-hcb, the exact multiplication of mu, and the displayed d bounds gives ||gamma\_jl,n(XY)-gamma\_jk,n(X) dot gamma\_kl,n(Y)||<=5(C\_H+kappa)e||X||||Y||. The unit and upper/lower norm clauses of conj-hcb similarly give diagonal-unit error <=3(C\_H+kappa)e and (1-(C\_H+kappa)e)||X||<=||gamma\_jk,n(X)||<=(1+2(C\_H+kappa)e)||X||. Since ||P\_1+P\_2-I||<=delta<=e, all clauses of def-four-corner-merging-datum hold with one common rho=D\_0*e, D\_0=5(C\_H+kappa), after harmless enlargement. Details are supplied by children 1.5.1--1.5.3. \steptag{validated} \\[6pt]
\step{1.5.1} Comparison, fixed amplifications, and exact adjoints. For (j,k)!=(1,1), h\_jk,n gamma\_jk,n=mu\_jk,n because node 1.4 gives h\_jk,n\textasciicircum{}\{-1\}=I\_n tensor h\_jk\textasciicircum{}\{-1\}; for (1,1), node 1.3 gives ||h\_11,n v\_n-mu\_11,n||<=kappa*e. Thus the d bounds in node 1.5 hold at every n for the same four level-one maps. For the off-diagonal pair, exact H-CB adjointness gives h\_kj,n(gamma\_jk,n(X)\textasciicircum{}dagger)=h\_jk,n(gamma\_jk,n(X))\textasciicircum{}dagger=mu\_jk,n(X)\textasciicircum{}dagger=mu\_kj,n(X\textasciicircum{}dagger)=h\_kj,n gamma\_kj,n(X\textasciicircum{}dagger); injectivity of h\_kj,n gives exact equality of the arguments. The diagonal transported corner is identical, while gamma\_11=v preserves dagger by def-extended-delta-inclusion. Hence the involution clause is exact. \steptag{validated} \\[6pt]
\step{1.5.1.1} Dependency bridge and complete derivation, valid only after nodes 1.3 and 1.4 are validated on the full EXT-CB scope (including e=0). Node 1.3 then supplies one fixed level-one spatial system mu\_jk with mu\_jk,n=I\_n tensor mu\_jk and ||h\_11,n v\_n-mu\_11,n||<=kappa*e for every n; node 1.4 supplies level-one bijectivity of every h\_jk, bijectivity (hence injectivity) of every h\_jk,n, and h\_jk,n\textasciicircum{}\{-1\}=I\_n tensor h\_jk\textasciicircum{}\{-1\}. Define gamma\_11=v and gamma\_jk=h\_jk\textasciicircum{}\{-1\} mu\_jk otherwise. Hence gamma\_jk,n=I\_n tensor gamma\_jk and, for (j,k)!=(1,1), h\_jk,n gamma\_jk,n=(I\_n tensor h\_jk)(I\_n tensor h\_jk\textasciicircum{}\{-1\}mu\_jk)=mu\_jk,n, while d\_11,n=h\_11,n v\_n-mu\_11,n has norm at most kappa*e. Thus all comparison bounds hold also at e=0. For every (j,k)!=(1,1), including (2,2), exact H-CB adjointness and spatial adjoint compatibility give h\_kj,n(gamma\_jk,n(X)\textasciicircum{}dagger)=h\_jk,n(gamma\_jk,n(X))\textasciicircum{}dagger=mu\_jk,n(X)\textasciicircum{}dagger=mu\_kj,n(X\textasciicircum{}dagger)=h\_kj,n gamma\_kj,n(X\textasciicircum{}dagger); injectivity of h\_kj,n yields gamma\_jk,n(X)\textasciicircum{}dagger=gamma\_kj,n(X\textasciicircum{}dagger). For (1,1), gamma\_11,n=v\_n preserves dagger because every amplification of an extended delta-inclusion is a delta-homomorphism, which in the star-algebra setting preserves dagger exactly. Therefore the four fixed maps satisfy the exact amplified involution clause. \steptag{validated} \\[6pt]
\step{1.5.1.1.1} Corrected validated endpoint bridge. No endpoint conclusion is inferred from the strict comparison written in node 1.3. Instead use validated child 1.5.1.1.1.2, whose explicit validated prerequisites are nodes 1.3.1.1, 1.3.2, and 1.4. Put T=h\_11 composed with v. Node 1.3.1.1 supplies one fixed mu\_11 with ||mu\_11,n-T\_n||<=kappa*e for every e>=0 and every n: it takes mu\_11=T when e=0, while for e>0 its strict corrected estimate implies the non-strict bound. Node 1.3.2 supplies the same mu\_11 as one spatial four-corner system mu\_jk with fixed amplifications and exact adjoint compatibility. Node 1.4 supplies h\_jk,n\textasciicircum{}\{-1\}=I\_n tensor h\_jk\textasciicircum{}\{-1\} and injectivity. Define gamma\_11=v and gamma\_jk=h\_jk\textasciicircum{}\{-1\} composed with mu\_jk off (1,1). Then h\_jk,n gamma\_jk,n=mu\_jk,n off (1,1); d\_11,n=h\_11,n v\_n-mu\_11,n has norm at most kappa*e, and every other d\_jk,n is zero. For each off-(1,1) corner, exact H-CB adjointness and spatial adjoint compatibility give h\_kj,n(gamma\_jk,n(X)\textasciicircum{}dagger)=h\_kj,n gamma\_kj,n(X\textasciicircum{}dagger), so injectivity yields gamma\_jk,n(X)\textasciicircum{}dagger=gamma\_kj,n(X\textasciicircum{}dagger). For (1,1), gamma\_11,n=v\_n preserves dagger exactly. Thus the comparison and exact amplified involution conclusions hold on the full e>=0 scope, without using the impossible strict estimate at e=0. \steptag{validated} \\[6pt]
\step{1.5.1.1.1.1} Formal validation gate. Require validated nodes 1.3 and 1.4. Validated 1.3 includes the full endpoint split: if e=0 then delta=epsilon=0, the defects of T vanish, and mu\_11=T gives ||mu\_11,n-h\_11,n v\_n||=0=kappa*e; if e>0 the validated correction step gives a strict bound and therefore the required non-strict bound. Validated 1.4 gives h\_jk,n\textasciicircum{}\{-1\}=I\_n tensor h\_jk\textasciicircum{}\{-1\} and injectivity. Hence h\_jk,n gamma\_jk,n=mu\_jk,n off (1,1), d\_11,n=h\_11,n v\_n-mu\_11,n has norm at most kappa*e, and all other d\_jk,n vanish. Exact H-CB adjointness, spatial adjoint compatibility, and injectivity then yield gamma\_jk,n(X)\textasciicircum{}dagger=gamma\_kj,n(X\textasciicircum{}dagger) off (1,1), including (2,2); for (1,1), gamma\_11,n=v\_n preserves dagger exactly. This establishes the comparison and involution conclusions on every e>=0 without any strict inequality at the zero endpoint. \steptag{validated} \\[4pt]
\step{1.5.1.1.1.2} Corrected endpoint dependency gate. This step does not use the strict comparison written in node 1.3. Require validated nodes 1.3.1.1, 1.3.2, and 1.4. Put T=h\_11 composed with v. Validated node 1.3.1.1 supplies one level-one unital dagger-homomorphism mu\_11 and its fixed amplifications with ||mu\_11,n-T\_n||<=kappa*e for the full e>=0 scope: at e=0 it explicitly takes mu\_11=T and obtains norm 0=kappa*e, while at e>0 the validated correction gives a strict bound and hence this non-strict bound. Validated node 1.3.2 spatializes that same mu\_11 by one unitary U\_1 and supplies the four fixed maps mu\_jk with mu\_jk,n=I\_n tensor mu\_jk, exact spatial multiplication, and mu\_jk,n(X)\textasciicircum{}dagger=mu\_kj,n(X\textasciicircum{}dagger). Validated node 1.4 supplies h\_jk,n\textasciicircum{}\{-1\}=I\_n tensor h\_jk\textasciicircum{}\{-1\} and injectivity of every h\_jk,n. Define gamma\_11=v and gamma\_jk=h\_jk\textasciicircum{}\{-1\} composed with mu\_jk off (1,1). Then h\_jk,n gamma\_jk,n=mu\_jk,n off (1,1), while d\_11,n=h\_11,n v\_n-mu\_11,n has ||d\_11,n(X)||<=kappa*e||X||; thus all other d\_jk,n vanish. For every corner off (1,1), including (2,2), exact H-CB adjointness and the spatial adjoint identity give h\_kj,n(gamma\_jk,n(X)\textasciicircum{}dagger)=h\_jk,n(gamma\_jk,n(X))\textasciicircum{}dagger=mu\_jk,n(X)\textasciicircum{}dagger=mu\_kj,n(X\textasciicircum{}dagger)=h\_kj,n gamma\_kj,n(X\textasciicircum{}dagger), so injectivity of h\_kj,n yields gamma\_jk,n(X)\textasciicircum{}dagger=gamma\_kj,n(X\textasciicircum{}dagger). For (1,1), gamma\_11,n=v\_n preserves dagger exactly by def-extended-delta-inclusion. Therefore the comparison and exact amplified involution conclusions hold for every e>=0, and no endpoint conclusion is attributed to node 1.3. \steptag{validated} \\[4pt]
\step{1.5.1.1.1.3} Dependency-adoption gate addressing ch-598304ea60cdaa3e. The amended parent expressly discards the false reading of node 1.3 and derives its conclusion from already validated child 1.5.1.1.1.2. That child has validation\_deps [1.3.1.1,1.3.2,1.4], explicitly handles e=0 by mu\_11=T and 0=kappa*e, uses the positive-e correction only for e>0, spatializes that same mu\_11, and obtains all Ha inverses and injectivity from node 1.4. Therefore the amended parent now adopts the corrected dependency chain and makes no inference from node 1.3 strict comparison. \steptag{validated} \\[4pt]
\step{1.5.2} Product estimate with every error term displayed. First ||gamma\_jk,n(X)||<=||h\_jk,n\textasciicircum{}\{-1\}||*||mu\_jk,n(X)+d\_jk,n(X)||<=2||X|| when C\_H*e and kappa*e are at most 1/4. Put D=gamma\_jl,n(XY)-gamma\_jk,n(X) dot gamma\_kl,n(Y). The conj-hcb product defect bounds the difference between h\_jl,n of the second term and (h\_jk,n gamma\_jk,n(X))(h\_kl,n gamma\_kl,n(Y)) by 4C\_H*e||X||||Y||. Substituting h gamma=mu+d, using mu\_jl(XY)=mu\_jk(X)mu\_kl(Y), and ||d\_ab,n(Z)||<=kappa*e||Z|| gives ||h\_jl,n(D)|| <= [kappa+2kappa+kappa\textasciicircum{}2*e+4C\_H]e||X||||Y|| <=4(C\_H+kappa)e||X||||Y||. Node 1.4 gives ||h\_jl,n\textasciicircum{}\{-1\}||<=1+C\_H*e<=5/4, hence ||D||<=5(C\_H+kappa)e||X||||Y||, uniformly in every index and amplification. \steptag{validated} \\[6pt]
\step{1.5.2.1} Prerequisites, with validation dependency made explicit. Require nodes 1.3 and 1.4 to be validated. Node 1.3 supplies one fixed spatial system mu\_jk,n=I\_n tensor mu\_jk and the uniform comparison ||h\_11,n v\_n-mu\_11,n||<=kappa e, including its separate e=0 branch. Node 1.4 supplies, for every j,k,n, bijectivity of h\_jk,n, h\_jk,n\textasciicircum{}\{-1\}=I\_n tensor h\_jk\textasciicircum{}\{-1\}, and ||h\_jk,n\textasciicircum{}\{-1\}||<=1+C\_H e. For the four maps already defined in node 1.5, gamma\_11=v and gamma\_jk=h\_jk\textasciicircum{}\{-1\} mu\_jk off (1,1), define d\_jk,n=h\_jk,n gamma\_jk,n-mu\_jk,n. Then d\_11,n=h\_11,n v\_n-mu\_11,n, so ||d\_11,n(Z)||<=kappa e||Z||, while d\_jk,n=0 off (1,1), because h\_jk,n(I\_n tensor h\_jk\textasciicircum{}\{-1\})=I. Thus uniformly ||d\_jk,n(Z)||<=kappa e||Z|| and every inverse used below exists with the stated bound; none of these facts is imported from an unvalidated sibling. \steptag{validated} \\[6pt]
\step{1.5.2.1.1} Validation-gated derivation of the prerequisite package. This child has validation dependencies 1.3 and 1.4 and therefore supplies no conclusion while either node is pending. After both validate, node 1.3 supplies fixed level-one maps mu\_jk with mu\_jk,n=I\_n tensor mu\_jk and ||(h\_11,n v\_n-mu\_11,n)(X)||<=kappa e||X|| for every n and X, on the full e>=0 scope. Node 1.4 supplies bijective h\_jk,n, the identity h\_jk,n\textasciicircum{}\{-1\}=I\_n tensor h\_jk\textasciicircum{}\{-1\}, and ||h\_jk,n\textasciicircum{}\{-1\}||<=1+C\_H e. Define gamma\_11=v and, off (1,1), gamma\_jk=h\_jk\textasciicircum{}\{-1\} composed with mu\_jk. Then gamma\_jk,n=(I\_n tensor h\_jk\textasciicircum{}\{-1\}) composed with mu\_jk,n off (1,1), so h\_jk,n gamma\_jk,n=(I\_n tensor h\_jk)(I\_n tensor h\_jk\textasciicircum{}\{-1\})mu\_jk,n=mu\_jk,n. Therefore d\_jk,n:=h\_jk,n gamma\_jk,n-mu\_jk,n is zero off (1,1), whereas d\_11,n=h\_11,n v\_n-mu\_11,n and hence ||d\_11,n(X)||<=kappa e||X||. Consequently ||d\_jk,n(X)||<=kappa e||X|| for every corner and amplification, and every inverse used in node 1.5.2 exists with the stated uniform bound. The conclusion is deliberately unavailable until validation of 1.3, so no pending correction result is used. \steptag{validated} \\[4pt]
\step{1.5.2.2} Conditional norm calculation from the now-declared prerequisites. Fix compatible indices j,k,l, an amplification n, X in M\_n tensor B(K\_k,K\_j), and Y in M\_n tensor B(K\_l,K\_k). Shrink the universal threshold so C\_H e<=1/4 and kappa e<=1/4. By node 1.5.2.1, h\_jk,n gamma\_jk,n(X)=mu\_jk,n(X)+d\_jk,n(X), ||h\_jk,n\textasciicircum{}\{-1\}||<=1+C\_H e, and ||d\_jk,n(X)||<=kappa e||X||. Since every mu\_jk,n is isometric, ||gamma\_jk,n(X)||<= (1+C\_H e)(1+kappa e)||X||<=25/16||X||<=2||X||, and similarly for Y. Put D=gamma\_jl,n(XY)-gamma\_jk,n(X) dot gamma\_kl,n(Y). The product-defect clause of conj-hcb gives E=h\_jl,n(gamma\_jk,n(X) dot gamma\_kl,n(Y))-(h\_jk,n gamma\_jk,n(X))(h\_kl,n gamma\_kl,n(Y)) with ||E||<=C\_H e||gamma\_jk,n(X)|| ||gamma\_kl,n(Y)||<=4C\_H e||X||||Y||. Hence, using exact spatial multiplication mu\_jl,n(XY)=mu\_jk,n(X)mu\_kl,n(Y), h\_jl,n(D)=d\_jl,n(XY)-mu\_jk,n(X)d\_kl,n(Y)-d\_jk,n(X)mu\_kl,n(Y)-d\_jk,n(X)d\_kl,n(Y)-E. Submultiplicativity ||XY||<=||X||||Y||, isometry of mu, and the d-bounds yield ||h\_jl,n(D)||<=[3kappa+kappa\textasciicircum{}2 e+4C\_H]e||X||||Y||. Because kappa e<=1, kappa\textasciicircum{}2 e<=kappa, so this is <=4(C\_H+kappa)e||X||||Y||. Finally node 1.5.2.1 gives ||D||<=||h\_jl,n\textasciicircum{}\{-1\}||||h\_jl,n(D)||<=(5/4)4(C\_H+kappa)e||X||||Y||=5(C\_H+kappa)e||X||||Y||. Every constant and threshold is independent of j,k,l,n,r and dim A. \steptag{validated} \\[6pt]
\step{1.5.2.2.1} Dependency-gated product calculation (addressing ch-e5329953fda228a2). This child has validation dependency 1.5.2.1 and therefore yields no product estimate while 1.5.2.1 is pending. Once 1.5.2.1 is validated, it supplies, for every compatible j,k,l,n, the fixed spatial isometries mu\_ab,n, bijectivity with ||h\_ab,n\textasciicircum{}\{-1\}||<=1+C\_H e, and d\_ab,n:=h\_ab,n gamma\_ab,n-mu\_ab,n satisfying ||d\_ab,n(Z)||<=kappa e||Z||. Under C\_H e<=1/4 and kappa e<=1/4, ||gamma\_jk,n(X)||<=||h\_jk,n\textasciicircum{}\{-1\}||( ||mu\_jk,n(X)||+||d\_jk,n(X)|| )<=(25/16)||X||<=2||X||, and similarly ||gamma\_kl,n(Y)||<=(25/16)||Y||<=2||Y||. Define D=gamma\_jl,n(XY)-gamma\_jk,n(X) dot gamma\_kl,n(Y), and let E be the H-CB product defect, so ||E||<=4C\_H e||X||||Y||. Linearity of h\_jl,n, the identities h\_ab,n gamma\_ab,n=mu\_ab,n+d\_ab,n and mu\_jl,n(XY)=mu\_jk,n(X)mu\_kl,n(Y) give exactly h\_jl,n(D)=d\_jl,n(XY)-mu\_jk,n(X)d\_kl,n(Y)-d\_jk,n(X)mu\_kl,n(Y)-d\_jk,n(X)d\_kl,n(Y)-E. Exact operator submultiplicativity in the matrix source, complete isometry of mu, and the d-bounds imply ||h\_jl,n(D)||<=[3kappa+kappa\textasciicircum{}2 e+4C\_H]e||X||||Y||<=4(C\_H+kappa)e||X||||Y||, since kappa e<=1. Applying the supplied inverse bound and C\_H e<=1/4 gives ||D||<=5(C\_H+kappa)e||X||||Y||. Hence the mathematics is valid after, and only after, validation of 1.5.2.1; the dependency is not replaced by any external input. \steptag{validated} \\[6pt]
\step{1.5.2.2.1.1} Ledger-enforced prerequisite bridge. This child has validation dependency 1.5.2.1, so it cannot be accepted before 1.5.2.1 is validated. After that validation, node 1.5.2.1 supplies the fixed spatial complete isometries mu\_ab,n, the inverse bounds ||h\_ab,n\textasciicircum{}\{-1\}||<=1+C\_H e, and ||d\_ab,n(Z)||<=kappa e||Z|| for d\_ab,n=h\_ab,n gamma\_ab,n-mu\_ab,n. If C\_H e<=1/4 and kappa e<=1/4, then ||gamma\_ab,n(Z)||<2||Z||. For D=gamma\_jl,n(XY)-gamma\_jk,n(X) dot gamma\_kl,n(Y), the H-CB product defect E and exact spatial multiplication give h\_jl,n(D)=d\_jl,n(XY)-mu\_jk,n(X)d\_kl,n(Y)-d\_jk,n(X)mu\_kl,n(Y)-d\_jk,n(X)d\_kl,n(Y)-E. Hence ||h\_jl,n(D)||<=[3kappa+kappa\textasciicircum{}2 e+4C\_H]e||X||||Y||<=4(C\_H+kappa)e||X||||Y||, and applying ||h\_jl,n\textasciicircum{}\{-1\}||<=5/4 yields ||D||<=5(C\_H+kappa)e||X||||Y||. Thus the parent calculation follows, but only after the ledger-enforced prerequisite validates. \steptag{archived} \\[4pt]
\step{1.5.2.2.1.2} Actual validation gate. This child may be accepted only after node 1.5.2.1 is validated. Once that happens, 1.5.2.1 provides the fixed spatial complete isometries mu\_ab,n, the bounds ||h\_ab,n\textasciicircum{}\{-1\}||<=1+C\_H e, and ||d\_ab,n(Z)||<=kappa e||Z||. Under C\_H e<=1/4 and kappa e<=1/4, the corrected calculation is ||gamma\_ab,n(Z)||<=(25/16)||Z||<=2||Z|| for every Z; the H-CB defect identity gives ||h\_jl,n(D)||<=[3kappa+kappa\textasciicircum{}2 e+4C\_H]e||X||||Y||<=4(C\_H+kappa)e||X||||Y||; and ||h\_jl,n\textasciicircum{}\{-1\}||<=5/4 gives ||D||<=5(C\_H+kappa)e||X||||Y||. No spatial map, inverse estimate, d-bound, or product conclusion is asserted before that prerequisite validates. \steptag{validated} \\[6pt]
\step{1.5.2.2.1.2.1} For every Z, dependency 1.5.2.1 gives gamma\_ab,n(Z)=h\_ab,n\textasciicircum{}\{-1\}(mu\_ab,n(Z)+d\_ab,n(Z)), with ||h\_ab,n\textasciicircum{}\{-1\}||<=1+C\_H e<=5/4, ||mu\_ab,n(Z)||=||Z|| by complete isometry, and ||d\_ab,n(Z)||<=kappa e||Z||<=||Z||/4. Therefore ||gamma\_ab,n(Z)||<=(5/4)(1+1/4)||Z||=(25/16)||Z||<=2||Z||. This non-strict estimate includes Z=0 and is the bound used in the H-CB defect estimate; the remaining displayed arithmetic is unchanged. \steptag{archived} \\[4pt]
\step{1.5.2.2.1.2.2} For every Z, validated node 1.5.2.1 gives gamma\_ab,n(Z)=h\_ab,n\textasciicircum{}\{-1\}(mu\_ab,n(Z)+d\_ab,n(Z)), with ||h\_ab,n\textasciicircum{}\{-1\}||<=1+C\_H e<=5/4, ||mu\_ab,n(Z)||=||Z|| by complete isometry, and ||d\_ab,n(Z)||<=kappa e||Z||<=||Z||/4. Therefore ||gamma\_ab,n(Z)||<=(5/4)(1+1/4)||Z||=(25/16)||Z||<=2||Z||. This non-strict estimate includes Z=0 and is exactly sufficient for the H-CB defect bound; all remaining arithmetic in the amended parent follows unchanged. \steptag{validated} \\[4pt]
\step{1.5.2.2.1.3} Challenge-closing dependency gate. This proof step is ledger-blocked on validation of 1.5.2.1. Only after that validation do the spatial maps mu\_ab,n, the inverse estimates ||h\_ab,n\textasciicircum{}\{-1\}||<=1+C\_H e, and the uniform d\_ab,n bounds enter scope. With those premises, the parent identity and estimates give ||D||<=5(C\_H+kappa)e||X||||Y|| exactly as displayed; before validation, this child yields no product estimate. \steptag{archived} \\[4pt]
\step{1.5.3} Units, norms, bijectivity, and datum assembly. Let u\_j be the compressed unit in S\_\{P\_j\}. The unit clauses of conj-hcb give ||h\_jj,n(I\_n tensor u\_j)-I||<=C\_H*e and ||u\_j-P\_j||<=C\_H*e after enlarging C\_H. Since h\_jj,n gamma\_jj,n(I)=I+d\_jj,n(I), node 1.4 yields ||gamma\_jj,n(I)-I\_n tensor P\_j||<=3(C\_H+kappa)e. Also ||mu\_jk,n(X)+d\_jk,n(X)|| lies between (1-kappa*e)||X|| and (1+kappa*e)||X||. Combining this with the upper bound and inverse bound for h\_jk,n gives (1-(C\_H+kappa)e)||X||<=||gamma\_jk,n(X)||<=(1+2(C\_H+kappa)e)||X|| after universal threshold reduction. The map gamma\_11=v is bijective; every other gamma\_jk is a composition of the bijections mu\_jk and h\_jk\textasciicircum{}\{-1\}. Together with node 1.5.1, node 1.5.2, the delta-projection hypotheses, and ||P\_1+P\_2-I||<=delta<=D\_0*e, these are exactly def-four-corner-merging-datum with common rho=D\_0*e, D\_0=5(C\_H+kappa), for four fixed level-one maps at all amplifications. \steptag{validated} \\[6pt]
\step{1.5.3.1} Dependency bridge for units, norms, bijectivity, and fixed amplifications, using only validated nodes 1.3.1, 1.3.2, and 1.4. For the same spatial map mu\_11 fixed in nodes 1.3.1 and 1.3.2, node 1.3.1 gives ||h\_11,n v\_n(X)-mu\_11,n(X)||<=kappa*e||X|| for every n and X. Define gamma\_11=v and, for (j,k)!=(1,1), gamma\_jk=h\_jk\textasciicircum{}\{-1\} composed with mu\_jk. Nodes 1.3.2 and 1.4 give mu\_jk,n=I\_n tensor mu\_jk and h\_jk,n\textasciicircum{}\{-1\}=I\_n tensor h\_jk\textasciicircum{}\{-1\}; hence every gamma\_jk,n is the amplification of the same level-one gamma\_jk, and h\_jk,n gamma\_jk,n=mu\_jk,n off (1,1). Thus, writing d\_jk,n=h\_jk,n gamma\_jk,n-mu\_jk,n, one has d\_jk,n=0 off (1,1) and ||d\_11,n(X)||<=kappa*e||X||. For diagonal source identity I, spatiality gives mu\_jj,n(I)=I. Consequently ||h\_jj,n(gamma\_jj,n(I)-I\_n tensor u\_j)||<=||d\_jj,n(I)||+||I-h\_jj,n(I\_n tensor u\_j)||<=(kappa+C\_H)e. Applying ||h\_jj,n\textasciicircum{}\{-1\}||<=1+C\_H*e and adding ||I\_n tensor(u\_j-P\_j)||<=C\_H*e yields ||gamma\_jj,n(I)-I\_n tensor P\_j||<=3(C\_H+kappa)e after shrinking the universal threshold so C\_H*e<=1. For arbitrary X, complete isometry of mu and the d-bound give (1-kappa*e)||X||<=||h\_jk,n gamma\_jk,n(X)||<=(1+kappa*e)||X||. The upper and inverse norm bounds for h from node 1.4 then give (1-(C\_H+kappa)e)||X||<=||gamma\_jk,n(X)||<=(1+2(C\_H+kappa)e)||X|| after a universal threshold reduction. Finally gamma\_11=v is bijective by the EXT-CB hypothesis, while every off-(1,1) gamma\_jk is the composition of the bijections h\_jk\textasciicircum{}\{-1\} and mu\_jk. Therefore the comparison, unit and norm estimates, bijectivity, and fixed-amplification conclusions follow without node 1.5.1. \steptag{archived} \\[6pt]
\step{1.5.3.1.1} Corrected availability statement (challenge ch-5031992bb44865d2). The missing comparison d\_11,n=h\_11,n v\_n-mu\_11,n with ||d\_11,n(X)||<=kappa*e*||X||, supplied by node 1.5.1 and ultimately node 1.3, prevents use of the displayed d\_11-based (1,1) unit/norm comparison in node 1.5.3.1 until those nodes are validated. It does not prevent bijectivity or fixed-amplification conclusions. Indeed gamma\_11=v is bijective by the root hypothesis and gamma\_11,n=I\_n tensor v by definition of amplification. For (j,k)!=(1,1), node 1.3.2 gives a fixed bijection mu\_jk with mu\_jk,n=I\_n tensor mu\_jk, and node 1.4 gives a fixed bijection h\_jk with h\_jk,n\textasciicircum{}\{-1\}=I\_n tensor h\_jk\textasciicircum{}\{-1\}. Hence gamma\_jk=h\_jk\textasciicircum{}\{-1\} composed with mu\_jk is bijective and gamma\_jk,n=h\_jk,n\textasciicircum{}\{-1\} composed with mu\_jk,n=I\_n tensor gamma\_jk, independently of node 1.5.1 and d\_11. After nodes 1.3 and 1.5.1 are validated, their d\_11 comparison also licenses the conditional triangle-inequality calculation in node 1.5.3.1. No permitted external input supplies that missing comparison, and no independently established bijectivity or fixed-amplification conclusion is withheld because of it. \steptag{validated} \\[6pt]
\step{1.5.3.1.1.1} Explicit bridge separating the available algebraic conclusions from the unavailable comparison estimate. For gamma\_11=v, the EXT-CB root hypothesis makes gamma\_11 bijective and amplification means gamma\_11,n=I\_n tensor gamma\_11. For (j,k)!=(1,1), validated node 1.3.2 gives bijective mu\_jk with mu\_jk,n=I\_n tensor mu\_jk, while validated node 1.4 gives bijective h\_jk and h\_jk,n\textasciicircum{}\{-1\}=I\_n tensor h\_jk\textasciicircum{}\{-1\}. Therefore gamma\_jk=h\_jk\textasciicircum{}\{-1\} mu\_jk is bijective and I\_n tensor gamma\_jk=(I\_n tensor h\_jk\textasciicircum{}\{-1\})(I\_n tensor mu\_jk)=h\_jk,n\textasciicircum{}\{-1\} mu\_jk,n=gamma\_jk,n. None of these equalities uses d\_11 or node 1.5.1. By contrast, the particular (1,1) comparison h\_11,n gamma\_11,n=mu\_11,n+d\_11,n with ||d\_11,n||<=kappa*e is exactly the pending content of node 1.5.1 (ultimately node 1.3), so only arguments invoking that comparison must wait. This removes the former invalid inference that all four kinds of conclusion were blocked. \steptag{validated} \\[4pt]
\step{1.5.3.1.2} Replacement validated-premise bridge addressing ch-050f3881bf97e920. Validated node 1.3.1 supplies the same level-one map mu\_11 used in validated node 1.3.2 and, for every n and X, ||h\_11,n v\_n(X)-mu\_11,n(X)||<=kappa*e||X||. Define gamma\_11=v and, for (j,k)!=(1,1), gamma\_jk=h\_jk\textasciicircum{}\{-1\} composed with mu\_jk. Validated node 1.4 gives h\_jk,n\textasciicircum{}\{-1\}=I\_n tensor h\_jk\textasciicircum{}\{-1\}; hence gamma\_jk,n=I\_n tensor gamma\_jk and h\_jk,n gamma\_jk,n=mu\_jk,n off (1,1). Thus d\_jk,n:=h\_jk,n gamma\_jk,n-mu\_jk,n vanishes off (1,1), while d\_11,n=h\_11,n v\_n-mu\_11,n and ||d\_11,n(X)||<=kappa*e||X||. This derives the formerly unavailable comparison without node 1.5.1. For diagonal source identity I, spatiality gives mu\_jj,n(I)=I, so linearity and the H-CB unit estimate imply ||h\_jj,n(gamma\_jj,n(I)-I\_n tensor u\_j)||<=||d\_jj,n(I)||+||I-h\_jj,n(I\_n tensor u\_j)||<=(kappa+C\_H)e. Applying ||h\_jj,n\textasciicircum{}\{-1\}||<=1+C\_H*e and adding ||I\_n tensor(u\_j-P\_j)||<=C\_H*e yields ||gamma\_jj,n(I)-I\_n tensor P\_j||<=3(C\_H+kappa)e once C\_H*e<=1. For arbitrary X, complete isometry of mu and the d-bound give (1-kappa*e)||X||<=||h\_jk,n gamma\_jk,n(X)||<=(1+kappa*e)||X||. From ||h\_jk,n||<=1+C\_H*e and ||h\_jk,n\textasciicircum{}\{-1\}||<=1+C\_H*e, therefore ||gamma\_jk,n(X)||>=(1-kappa*e)/(1+C\_H*e)||X||>=(1-(C\_H+kappa)e)||X|| and ||gamma\_jk,n(X)||<=(1+C\_H*e)(1+kappa*e)||X||<=(1+2(C\_H+kappa)e)||X|| after shrinking so C\_H*e<=1. Finally gamma\_11=v is bijective by the EXT-CB hypothesis; off (1,1), gamma\_jk is a composition of the validated bijections h\_jk\textasciicircum{}\{-1\} and mu\_jk. The displayed tensor identities prove all four are fixed level-one maps at every amplification. Hence all conclusions of node 1.5.3.1 follow from validated nodes 1.3.1, 1.3.2, and 1.4, independently of pending node 1.5.1. \steptag{validated} \\[6pt]
\step{1.5.3.1.2.1} Ledger-licensed modus-ponens bridge. Require validated nodes 1.3.1, 1.3.2, and 1.4. Node 1.3.1 supplies, for the same spatial map mu\_11 of node 1.3.2, ||h\_11,n v\_n(X)-mu\_11,n(X)|| <= kappa*e||X|| for every n and X. Set gamma\_11=v and gamma\_jk=h\_jk\textasciicircum{}\{-1\} composed with mu\_jk off (1,1). By node 1.4, h\_jk,n\textasciicircum{}\{-1\}=I\_n tensor h\_jk\textasciicircum{}\{-1\}; by node 1.3.2, mu\_jk,n=I\_n tensor mu\_jk. Hence gamma\_jk,n=I\_n tensor gamma\_jk and h\_jk,n gamma\_jk,n=mu\_jk,n off (1,1). Therefore d\_jk,n:=h\_jk,n gamma\_jk,n-mu\_jk,n is zero off (1,1), while d\_11,n=h\_11,n v\_n-mu\_11,n and ||d\_11,n(X)|| <= kappa*e||X||. For diagonal source identity I, spatiality gives mu\_jj,n(I)=I; the triangle inequality, the H-CB unit estimate, and ||h\_jj,n\textasciicircum{}\{-1\}|| <= 1+C\_H*e give ||gamma\_jj,n(I)-I\_n tensor u\_j|| <= (1+C\_H*e)(C\_H+kappa)e. Adding ||I\_n tensor(u\_j-P\_j)|| <= C\_H*e and shrinking the universal threshold so C\_H*e <= 1 gives ||gamma\_jj,n(I)-I\_n tensor P\_j|| <= 3(C\_H+kappa)e. For arbitrary X, complete isometry of mu and the d-bound give (1-kappa*e)||X|| <= ||h\_jk,n gamma\_jk,n(X)|| <= (1+kappa*e)||X||. Using ||h\_jk,n|| <= 1+C\_H*e and ||h\_jk,n\textasciicircum{}\{-1\}|| <= 1+C\_H*e yields (1-(C\_H+kappa)e)||X|| <= ||gamma\_jk,n(X)|| <= (1+2(C\_H+kappa)e)||X|| after the same universal threshold reduction. Finally gamma\_11=v is bijective by the EXT-CB hypothesis, and off (1,1) gamma\_jk is the composition of the bijections h\_jk\textasciicircum{}\{-1\} and mu\_jk. Thus the comparison, unit and norm estimates, bijectivity, and fixed-amplification conclusions follow solely from the three registered validated premises. \steptag{validated} \\[4pt]
\step{1.5.3.1.3} Ledger-enforced replacement for the pending 1.5.1 dependency. Require validated nodes 1.3.1, 1.3.2, and 1.4. Node 1.3.1 gives ||h\_11,n v\_n(X)-mu\_11,n(X)||<=kappa*e||X|| for every n and X for the same mu\_11 spatialized in node 1.3.2. With gamma\_11=v and gamma\_jk=h\_jk\textasciicircum{}\{-1\}mu\_jk off (1,1), node 1.4 gives h\_jk,n gamma\_jk,n=mu\_jk,n off (1,1); hence d\_11,n=h\_11,n v\_n-mu\_11,n and d\_jk,n=0 otherwise. Therefore ||d\_jk,n(X)||<=kappa*e||X|| without using node 1.5.1. The H-CB diagonal unit estimate, spatial identity mu\_jj,n(I)=I, and ||h\_jj,n\textasciicircum{}\{-1\}||<=1+C\_H*e then give ||gamma\_jj,n(I)-I\_n tensor P\_j||<=3(C\_H+kappa)e when C\_H*e<=1. Complete isometry of mu and the d-bound give (1-kappa*e)||X||<=||h\_jk,n gamma\_jk,n(X)||<=(1+kappa*e)||X||; the upper and inverse bounds for h yield (1-(C\_H+kappa)e)||X||<=||gamma\_jk,n(X)||<=(1+2(C\_H+kappa)e)||X|| under the same universal threshold reduction. Bijectivity and fixed amplification follow because gamma\_11=v and, off (1,1), gamma\_jk=h\_jk\textasciicircum{}\{-1\}mu\_jk, with h\_jk,n\textasciicircum{}\{-1\}=I\_n tensor h\_jk\textasciicircum{}\{-1\} and mu\_jk,n=I\_n tensor mu\_jk. Thus every conclusion of node 1.5.3.1 is licensed by these three validated prerequisites, and pending node 1.5.1 is unnecessary for this bridge. \steptag{validated} \\[4pt]
\step{1.5.3.2} Final datum assembly with explicit prerequisites. Require validated node 1.5.3.1 for bijectivity, fixed amplifications, diagonal-unit error, and two-sided norm control; validated node 1.5.1 for exact involution compatibility and the d-comparison; and validated node 1.5.2 for compressed-product defect at most 5(C\_H+kappa)e. The target elements P\_1=P and P\_2=Q are delta-projections by the root hypotheses and ||P\_1+P\_2-I||<=delta<=e<=D\_0*e, where D\_0=5(C\_H+kappa) after enlarging universal constants so D\_0>=1. Thus, with rho=D\_0*e, every clause of the amended def-four-corner-merging-datum holds for the four fixed level-one maps gamma\_jk and all their amplifications. No conclusion of this child is available until all three declared dependencies are validated. \steptag{archived} \\[6pt]
\step{1.5.3.2.1} Replacement assembly from validated descendants, so no pending premise is used. Require validated nodes 1.3.2, 1.4, 1.5.2.2.1.2, and 1.5.3.1.3. Node 1.5.3.1.3 supplies four bijective fixed level-one maps gamma\_jk, their fixed amplifications gamma\_jk,n=I\_n tensor gamma\_jk, diagonal-unit error at most 3(C\_H+kappa)e, and two-sided bounds (1-(C\_H+kappa)e)||X|| <= ||gamma\_jk,n(X)|| <= (1+2(C\_H+kappa)e)||X||. Node 1.5.2.2.1.2 supplies compressed-product defect at most 5(C\_H+kappa)e. It remains only to establish the exact involution clause without pending node 1.5.1. For gamma\_11=v, exact dagger preservation follows from def-extended-delta-inclusion. If (j,k)!=(1,1), exact H-CB adjointness, spatial adjoint compatibility from node 1.3.2, the identities h\_jk,n gamma\_jk,n=mu\_jk,n and h\_kj,n gamma\_kj,n=mu\_kj,n from node 1.5.3.1.3, and injectivity of h\_kj,n from node 1.4 give h\_kj,n(gamma\_jk,n(X)\textasciicircum{}dagger)=h\_jk,n(gamma\_jk,n(X))\textasciicircum{}dagger=mu\_jk,n(X)\textasciicircum{}dagger=mu\_kj,n(X\textasciicircum{}dagger)=h\_kj,n gamma\_kj,n(X\textasciicircum{}dagger), hence gamma\_jk,n(X)\textasciicircum{}dagger=gamma\_kj,n(X\textasciicircum{}dagger). This includes the (2,2) corner. Let D\_0=5(C\_H+kappa), enlarging the universal constants if needed so D\_0>=1, and rho=D\_0 e. The canonical source block projections in M\_\{r+1\}=B(K\_1 direct-sum K\_2) are exactly complementary. The targets P\_1=P and P\_2=Q are delta-projections by the root hypotheses, and ||P\_1+P\_2-I||<=delta<=e<=rho. Exact involution has defect 0<=rho; the product, unit, lower-norm, and upper-norm defects above are respectively bounded by D\_0 e, 3(C\_H+kappa)e, (C\_H+kappa)e, and 2(C\_H+kappa)e, all at most rho. Consequently every clause of the amended def-four-corner-merging-datum holds with the common defect rho for the same four level-one maps at every amplification. This is a complete replacement route and does not invoke pending nodes 1.5.3.1, 1.5.1, or 1.5.2. \steptag{validated} \\[4pt]
\step{1.5.3.3} Replacement dependency bridge for units, norms, bijectivity, and fixed amplifications. Require validated nodes 1.3.1, 1.3.2, and 1.4. For the same spatial map mu\_11 fixed in nodes 1.3.1 and 1.3.2, node 1.3.1 gives ||h\_11,n v\_n(X)-mu\_11,n(X)||<=kappa*e||X|| for every n and X. Define gamma\_11=v and, for (j,k)!=(1,1), gamma\_jk=h\_jk\textasciicircum{}\{-1\} composed with mu\_jk. Nodes 1.3.2 and 1.4 give mu\_jk,n=I\_n tensor mu\_jk and h\_jk,n\textasciicircum{}\{-1\}=I\_n tensor h\_jk\textasciicircum{}\{-1\}, so all gamma\_jk,n are fixed amplifications, h\_jk,n gamma\_jk,n=mu\_jk,n off (1,1), and d\_jk,n:=h\_jk,n gamma\_jk,n-mu\_jk,n is zero off (1,1) while ||d\_11,n(X)||<=kappa*e||X||. For diagonal identity I, mu\_jj,n(I)=I, hence ||h\_jj,n(gamma\_jj,n(I)-I\_n tensor u\_j)||<=||d\_jj,n(I)||+||I-h\_jj,n(I\_n tensor u\_j)||<=(kappa+C\_H)e. The inverse bound and ||u\_j-P\_j||<=C\_H*e imply ||gamma\_jj,n(I)-I\_n tensor P\_j||<=3(C\_H+kappa)e once C\_H*e<=1. For arbitrary X, complete isometry of mu gives (1-kappa*e)||X||<=||h\_jk,n gamma\_jk,n(X)||<=(1+kappa*e)||X||. The upper and inverse bounds for h then imply (1-(C\_H+kappa)e)||X||<=||gamma\_jk,n(X)||<=(1+2(C\_H+kappa)e)||X|| after a universal threshold reduction. Finally gamma\_11=v is bijective by the EXT-CB hypothesis, and every off-(1,1) gamma\_jk is the composition of bijections h\_jk\textasciicircum{}\{-1\} and mu\_jk. Thus all bridge conclusions follow without node 1.5.1. \steptag{validated} \\[4pt]
\step{1.5.3.4} Active adoption of the validated replacement assembly. Validated node 1.5.3.2.1 proves, without invoking pending nodes 1.5.3.1, 1.5.1, or 1.5.2, that the same four bijective fixed level-one maps gamma\_jk and their amplifications satisfy every clause of the amended def-four-corner-merging-datum with common defect rho=D\_0 e, where D\_0=5(C\_H+kappa) is universal and enlarged so D\_0>=1. Therefore the final datum-assembly conclusion required in node 1.5.3 follows directly from 1.5.3.2.1; this sibling is the active route and the stale node 1.5.3.2 is retired. \steptag{validated} \\[4pt]
\step{1.6} Merge and conclude the root. Choose e\_ext as the positive minimum of the thresholds in nodes 1.1--1.5, a\_merge/(D\_0+1), and the finitely many Neumann and 1/4 bounds; it is universal. Node 1.5 supplies four fixed bijective level-one corner maps satisfying def-four-corner-merging-datum with rho=D\_0*e and rho+epsilon<=(D\_0+1)e<=a\_merge. Apply lem-extcb-four-corner-merge (whose assembled norm component is lem-extcb-four-corner-norm) to obtain the single sum map v\_+:M\_\{r+1\}->A. It is an extended C\_merge*(rho+epsilon)-isomorphism, hence an extended C\_ext*e-isomorphism for C\_ext=C\_merge*(D\_0+1). Its n-th map is I\_n tensor v\_+, assembled from I\_n tensor U\_1, I\_n tensor U\_2 and the same four gamma\_jk. All constants used are universal and therefore independent of r,n,dim A and block data. This is exactly node 1. \steptag{validated} \\[4pt]
\end{proofbox}

\end{document}
